\chapter{El método de Malliavin--Stein}\label{cap:malliavin-stein}

\begin{panorama}
En 2005, Nualart y Peccati~\cite{NualartPeccati05} hallaron un fenómeno sorprendente: una sucesión de
variables en un caos de Wiener fijo, normalizada, converge a la normal \emph{si y solo si su
cuarto momento converge a} 3. Poco después, Nourdin y Peccati (2009) mostraron que
combinando el cálculo de Malliavin con el método de Stein se obtienen cotas \emph{cuantitativas}
explícitas. El ``método de Malliavin--Stein'' es hoy una de las áreas más activas del cálculo
de Malliavin, con aplicaciones a la estadística de procesos de memoria larga, la geometría
aleatoria, las EDPEs (Capítulo~\ref{cap:anderson}), las ondas aleatorias y las redes neuronales anchas.
\end{panorama}

\section{El método de Stein}

\begin{lema}[Stein]\label{lem:stein-ecuacion}
Sea $N\sim N(0,1)$ y $h:\R\to\R$ medible acotada. La ecuación de Stein
\[
  f'(x)-xf(x)=h(x)-\E h(N)
\]
tiene una solución acotada $f_h(x)=e^{x^2/2}\int_{-\infty}^x(h(y)-\E h(N))e^{-y^2/2}\,\dd y$. Si $h=\ind_A$ con $A$
boreliano, entonces $\norm{f_h}_\infty\le\sqrt{\pi/2}$ y $\norm{f_h'}_\infty\le2$. Si $h$ es 1-Lipschitz, $\norm{f_h'}_\infty\le\sqrt{2/\pi}$ y
$\norm{f_h''}_\infty\le2$. En consecuencia, para toda variable $X$,
\[
  d_{TV}(X,N)=\sup_A\abs{\P(X\in A)-\P(N\in A)}\le\sup_{f\in\F_{TV}}\abs{\E[f'(X)-Xf(X)]},
\]
con $\F_{TV}=\{f\in C^1:\ \norm f_\infty\le\sqrt{\pi/2},\ \norm{f'}_\infty\le2\}$, y análogamente para la distancia de
Wasserstein $d_W$, con $\F_W=\{\norm{f'}_\infty\le\sqrt{2/\pi}\}$.
\end{lema}

\begin{proof}[Esbozo]
Que $f_h$ resuelve la ecuación es una verificación directa: $(e^{-x^2/2}f_h)'=(h-\E h(N))e^{-x^2/2}$.
Las cotas se obtienen con las estimaciones de Mills, $\int_x^\infty e^{-y^2/2}\,\dd y\le\min(\sqrt{\pi/2},1/x)e^{-x^2/2}$,
usando que $\int_\R(h-\E h(N))e^{-y^2/2}=0$ para escribir $f_h$ con la cola izquierda o la derecha según
el signo de $x$. Para $h$ indicador no continuo, $f_h$ es solo Lipschitz; se aproxima. Finalmente,
evaluando la ecuación en $X$ y tomando esperanzas, $\P(X\in A)-\P(N\in A)=\E[f_h'(X)-
Xf_h(X)]$.
\end{proof}

\begin{intuicion}
El método de Stein convierte la pregunta ``¿está $X$ cerca de $N$?'' en ``¿satisface $X$
aproximadamente la identidad $\E[Xf(X)]=\E[f'(X)]$?''. Y el cálculo de Malliavin es
precisamente una máquina para calcular $\E[Xf(X)]$: por la representación de covarianzas,
\[
  \E[Xf(X)]=\E\big[f'(X)\ip{DX}{-DL^{-1}X}_\Hf\big].
\]
Así, $X$ es casi normal si $\ip{DX}{-DL^{-1}X}\approx1$. Los dos métodos encajan como una llave
en su cerradura.
\end{intuicion}

\section{La cota fundamental}

\begin{teorema}[Nourdin--Peccati, 2009]\label{teo:nourdin-peccati}
Sea $F\in\D^{1,2}$ con $\E F=0$ y $\E F^2=1$. Entonces
\[
  d_{TV}(F,N)\le2\,\E\Big|1-\ip{DF}{-DL^{-1}F}_\Hf\Big|\le2\sqrt{\Var\big(\ip{DF}{-DL^{-1}F}_\Hf\big)}.
\]
\end{teorema}

\begin{proof}
Para $f\in\F_{TV}$, por la Proposición~\ref{prop:covarianzas} y la regla de la cadena,
\[
  \E[Ff(F)]=\E\ip{Df(F)}{-DL^{-1}F}=\E\big[f'(F)\ip{DF}{-DL^{-1}F}\big].
\]
Así $\abs{\E[f'(F)-Ff(F)]}=\abs{\E[f'(F)(1-\ip{DF}{-DL^{-1}F})]}\le2\E\abs{1-\ip{DF}{-DL^{-1}F}}$. La segunda
desigualdad es Cauchy--Schwarz, pues $\E\ip{DF}{-DL^{-1}F}=\E F^2=1$. (Estrictamente, la
regla de la cadena se aplica a $f$ Lipschitz con la versión de la Proposición~\ref{prop:cadena}, lo que requiere
que $F$ tenga densidad, o un argumento de aproximación.)
\end{proof}

\section{El teorema del cuarto momento}
Si $F=I_q(f)$, entonces $L^{-1}F=-\frac1qF$, y por tanto $\ip{DF}{-DL^{-1}F}=\frac1q\norm{DF}^2_\Hf$.

\begin{lema}\label{lem:cuarto-momento}
Sea $F=I_q(f)$ con $f\in\Hf^{\odot q}$, $q\ge2$, y $\E F^2=q!\norm f^2=1$. Entonces
\begin{align*}
  \Var\Big(\frac1q\norm{DF}^2_\Hf\Big)&=\frac1{q^2}\sum_{r=1}^{q-1}r^2(r!)^2\binom qr^4(2q-2r)!\,\norm{f\tilde\otimes_rf}^2,\\
  \E F^4-3&=\frac3q\sum_{r=1}^{q-1}r\,(r!)^2\binom qr^4(2q-2r)!\,\norm{f\tilde\otimes_rf}^2.
\end{align*}
En consecuencia,
\[
  \Var\Big(\frac1q\norm{DF}^2_\Hf\Big)\le\frac{q-1}{3q}\big(\E F^4-3\big).
\]
\end{lema}

\begin{proof}[Esbozo]
Por el Teorema~\ref{teo:derivada-caos}, $D_tF=qI_{q-1}(f(\cdot,t))$. Por la fórmula del producto (Teorema~\ref{teo:producto})
aplicada a $I_{q-1}(f(\cdot,t))^2$ e integrando en $t$ (lo que añade una contracción),
\[
  \frac1q\norm{DF}^2=q\sum_{r=1}^q(r-1)!\binom{q-1}{r-1}^2I_{2q-2r}(f\tilde\otimes_rf).
\]
El término $r=q$ es la constante $q!\norm f^2=1$, y la varianza se calcula con la isometría;
reescribiendo $q(r-1)!\binom{q-1}{r-1}^2=\frac rq\cdot r!\binom qr^2$ se obtiene la primera fórmula. Para el cuarto momento,
$\E F^4=\E[F\cdot F^3]=3\E[F^2\cdot\frac1q\norm{DF}^2]$ (dualidad y regla de la cadena: $\E[F\cdot F^3]=
\E\ip{D(F^3)}{-DL^{-1}F}=3\E[F^2\frac1q\norm{DF}^2]$), y se desarrolla $F^2$ con la fórmula del producto; el
resultado es la segunda fórmula. Comparando término a término, el cociente entre el coeficiente
de $\norm{f\tilde\otimes_rf}^2$ en la varianza y en $\E F^4-3$ es $\frac{r^2/q^2}{3r/q}=\frac r{3q}\le\frac{q-1}{3q}$, lo que da la desigualdad.
\end{proof}

\begin{teorema}[Teorema del cuarto momento]\label{teo:cuarto-momento}
Sea $q\ge2$ y $F_n=I_q(f_n)$ con $\E F_n^2\to1$. Son
equivalentes:
\begin{enumerate}[label=(\roman*)]
  \item $F_n\tolaw N(0,1)$;
  \item $\E F_n^4\to3$;
  \item $\norm{f_n\otimes_rf_n}_{\Hf^{\otimes(2q-2r)}}\to0$ para $r=1,\dots,q-1$;
  \item $\norm{DF_n}^2_\Hf\to q$ en $L^2(\Omega)$.
\end{enumerate}
Además, si $\E F^2=1$, se tiene la cota cuantitativa
\[
  d_{TV}(F,N)\le2\sqrt{\frac{q-1}{3q}\big(\E F^4-3\big)}.
\]
\end{teorema}

\begin{proof}
La cota es el Teorema~\ref{teo:nourdin-peccati} combinado con el Lema~\ref{lem:cuarto-momento}; da (ii)$\Rightarrow$(i). (i)$\Rightarrow$(ii): por
el Teorema~\ref{teo:normas-caos} (hipercontractividad), $\sup_n\E\abs{F_n}^p<\infty$ para todo $p$, luego hay integrabilidad
uniforme de $F_n^4$ y $\E F_n^4\to\E N^4=3$. (ii)$\Leftrightarrow$(iii): por la fórmula del cuarto cumulante del
Lema~\ref{lem:cuarto-momento}, y porque $\norm{f\tilde\otimes_rf}$ y $\norm{f\otimes_rf}$ son comparables (identidad combinatoria). (iii)$\Leftrightarrow$(iv):
por la fórmula de la varianza (véase también~\cite{NualartOrtizLatorre08}).
\end{proof}

\begin{intuicion}
En el caos $q\ge2$, el cuarto cumulante $\kappa_4(F)=\E F^4-3(\E F^2)^2$ es siempre $\ge0$ (Lema~\ref{lem:cuarto-momento}):
las variables de caos tienen colas más pesadas que la normal. La normal es el único
``mínimo'' posible, y alcanzarlo asintóticamente solo es posible si el núcleo $f_n$ se ``dispersa'',
es decir, si sus contracciones se anulan. Es un \emph{principio de rigidez}. Notablemente, la
condición (iv) dice que la ``varianza local'' $\frac1q\norm{DF}^2$ se vuelve determinista.
\end{intuicion}

\begin{teorema}[Peccati--Tudor~\cite{PeccatiTudor05}]
Sea $F_n=(F^1_n,\dots,F^d_n)$ con $F^i_n=I_{q_i}(f^i_n)$ y $\E[F^i_nF^j_n]\to
C_{ij}$. Entonces $F_n\tolaw N(0,C)$ si y solo si cada componente converge a su marginal
gaussiana.
\end{teorema}

Es decir: en el caos de Wiener, \emph{la convergencia componente a componente implica la
convergencia conjunta}. Esto falla para variables generales.

\section{Aplicación: variación cuadrática del browniano fraccionario}
Sea $B^H$ un fBm y $X_k=n^H(B^H_{(k+1)/n}-B^H_{k/n})$, $k=0,\dots,n-1$, que forman una sucesión
gaussiana estacionaria con $\E X_k^2=1$ y covarianza
\[
  \rho(k)=\tfrac12\big(\abs{k+1}^{2H}-2\abs k^{2H}+\abs{k-1}^{2H}\big)\sim H(2H-1)\abs k^{2H-2}.
\]
Sea $V_n=\frac1{\sqrt n}\sum_{k<n}(X_k^2-1)=\frac1{\sqrt n}\sum_kH_2(X_k)$, que vive en el segundo caos.

\begin{teorema}[Breuer--Major; Nourdin--Peccati]
Si $H<3/4$, $V_n\tolaw N(0,\sigma_H^2)$ con
$\sigma_H^2=2\sum_{k\in\Z}\rho(k)^2$, y
\[
  d_{TV}\Big(\frac{V_n}{\sqrt{\Var V_n}},N\Big)\le C\cdot
  \begin{cases}
    n^{-1/2}, & H\in(0,\frac58),\\
    n^{-1/2}\log^{3/2}n, & H=\frac58,\\
    n^{4H-3}, & H\in(\frac58,\frac34).
  \end{cases}
\]
Si $H>3/4$, $n^{1-2H}V_n$ converge a una variable de \textup{Rosenblatt} (no gaussiana, en el segundo
caos).
\end{teorema}

\begin{proof}[Esbozo]
$V_n=I_2(f_n)$ con $f_n=\frac1{\sqrt n}\sum_ke_k\otimes e_k$, donde $e_k\in\Hf$ con $\ip{e_k}{e_l}=\rho(k-l)$. Entonces
$\Var V_n=\frac2n\sum_{k,l<n}\rho(k-l)^2\to\sigma_H^2$, finito si y solo si $\sum\rho^2<\infty$, es decir, $H<3/4$. Además
\[
  \norm{f_n\otimes_1f_n}^2=\frac1{n^2}\sum_{i,j,k,l<n}\rho(i-j)\rho(j-k)\rho(k-l)\rho(l-i)\le\frac Cn\Big(\sum_{\abs k<n}\abs{\rho(k)}^{4/3}\Big)^3
\]
por la desigualdad de Young para convoluciones. Con $\rho(k)\sim k^{2H-2}$, la suma es acotada si
$H<5/8$, crece como $\log n$ si $H=5/8$ y como $n^{(8H-5)/3}$ si $H>5/8$. Por el Teorema~\ref{teo:cuarto-momento} (en
$q=2$, $\kappa_4=48\norm{f\otimes_1f}^2$ y la cota de distancia es $\propto\sqrt{\kappa_4}$) se obtienen las tasas. Para $H>3/4$,
la renormalización $n^{1-2H}V_n$ hace que $f_n$ converja en $\Hf^{\otimes2}$ a un núcleo no nulo, cuya integral
doble es la variable de Rosenblatt.
\end{proof}

Este resultado tiene una consecuencia estadística inmediata: el estimador de $H$ basado en
la variación cuadrática es asintóticamente normal si y solo si $H<3/4$.

\section{Refinamientos}

\paragraph{Tasas óptimas.} Nourdin y Peccati (2015) probaron que, para $F_n$ en un caos fijo, $d_{TV}(F_n,N)\asymp
\max\{\abs{\E F_n^3},\ \E F_n^4-3\}$: la cota $\sqrt{\kappa_4}$ del Teorema~\ref{teo:cuarto-momento} no es óptima, y la tasa correcta es lineal
en los cumulantes.

\paragraph{Entropía.} Nourdin, Peccati y Swan (2014) mostraron que el teorema del cuarto momento
vale también en \emph{entropía relativa}, una noción de convergencia más fuerte, y con ello en total
variación de las densidades.

\paragraph{Desigualdades de Poincaré de segundo orden.} Para $F$ general (no en un caos fijo),
Nourdin, Peccati y Reinert~\cite{NourdinPeccatiReinert09} probaron
\[
  d_{TV}(F,N)\le\frac{\sqrt{10}}{2\sigma^2}\Big(\E\norm{D^2F}^4_{\mathrm{op}}\Big)^{1/4}\Big(\E\norm{DF}^4_\Hf\Big)^{1/4},
\]
donde $\norm{D^2F}_{\mathrm{op}}$ es la norma de operador de la segunda derivada como forma bilineal en $\Hf$. Esta
cota solo requiere estimar $DF$ y $D^2F$, y es la herramienta estándar para funcionales no lineales
(EDPEs, redes neuronales con anchura grande, estadísticos de campos gaussianos).

\paragraph{Universalidad.} Nourdin, Peccati y Reinert (2010) probaron que el caos de Wiener es \emph{universal}:
si una sucesión de sumas homogéneas multilineales $Q_n(\mathbf G)=\sum f_n(i_1,\dots,i_q)G_{i_1}\cdots G_{i_q}$ en
gaussianas satisface un TCL, entonces lo satisface $Q_n(\mathbf X)$ para cualesquiera $X_i$ independientes
centradas de varianza unidad (bajo una condición de influencia pequeña). La prueba combina
el teorema del cuarto momento con un principio de Lindeberg.

\paragraph{Más allá de la normal.} El método se extiende a aproximaciones por la ley Gamma (segundo
caos), por leyes en el segundo caos en general (Azmoodeh, Peccati, Poly), y al marco de \emph{triples
de Markov} (Ledoux; Azmoodeh, Campese, Poly), donde el operador $L$ es el generador de un
semigrupo de difusión con espectro discreto (por ejemplo, los polinomios de Jacobi o Laguerre),
lo que produce teoremas del cuarto momento para ``caos'' de otras medidas. Existe también un
teorema del cuarto momento en \emph{probabilidad libre} (Kemp, Nourdin, Peccati, Speicher, 2012),
con la ley semicircular como límite. Véase la reseña~\cite{AzmoodehPeccatiYang21}.

\paragraph{Geometría aleatoria.} El método de Malliavin--Stein ha sido decisivo en el estudio de los
conjuntos nodales de ondas aleatorias (longitud de las líneas nodales de las ondas planas de
Berry, autofunciones aleatorias en la esfera y el toro): Marinucci, Peccati, Rossi y Wigman
obtuvieron TCLs y descubrieron la \emph{cancelación de Berry}, en la que la proyección en el segundo
caos se anula y el comportamiento lo dicta el cuarto caos.
